\documentclass[a4paper,12pt]{letter}

\usepackage{amsmath}
\usepackage{amssymb}
\usepackage{nopageno}
\usepackage{amsmath}
\usepackage{enumitem}
% \usepackage{mathtools}
%\usepackage{eqnarray,amsmath}
\usepackage[a4paper, total={7.5in, 11.25in}]{geometry}

\begin{document}

Name: \hrulefill\ Neptun code: \hrulefill

\bigskip

\begin{center}
(KTXFI2EBNF) Physics II. Exam~4~Solutions\\
February~2,~2026
\end{center}

\bigskip
 
\begin{enumerate}[label=\Roman*.]
   
\item Moving Reference Frames. Special Theory of Relativity

\begin{enumerate}[label=\arabic*.]
\setcounter{enumii}{0} % Számláló beállítása

\item State the Galilean principle of relativity. \hfill (5~points)

\smallskip
\textit{Lecture on October 13, 2025, Slides 8--16.}

\item Write down Einstein’s relation expressing mass–energy equivalence. What is its physical meaning? \hfill (5~points)

\smallskip
\textit{Lecture on October 13, 2025, Slides 23--24.}

\end{enumerate}

% ------------------------------------------------------------------------------------------------------------------------------------

\item Limits of the Classical Conceptual Framework

\begin{enumerate}[label=\arabic*.]
\setcounter{enumii}{2} % Számláló beállítása

\item What is the photoelectric effect? Explain the mechanism of the photoelectric effect. \hfill (5~points)

\smallskip
\textit{Lecture on October 6, 2025, Slides 23--27.}
    
\item Describe the Compton experiment and provide a qualitative but precise explanation of the phenomenon of Compton scattering. \hfill (5~points)

\smallskip
\textit{Lecture on October 6, 2025, Slides 29--31.}
  
\item Formulate the de Broglie principle in your own words. \hfill (5~points)

\smallskip
\textit{Lecture on October 20, 2025, Slide 23.}


\end{enumerate}
  
% ------------------------------------------------------------------------------------------------------------------------------------

\item Elements of Quantum Mechanics. Physics of Condensed Matter. Theory of Electrical Conduction

\begin{enumerate}[label=\arabic*.]
\setcounter{enumii}{5} % Számláló beállítása

\item What is the tunneling effect? Describe the phenomenon in your own words. \hfill (5~points)

\smallskip
\textit{Lecture on November 10, 2025, Slide 46.}

\item What is a forbidden band (band gap)? Define the concept. \hfill (5~points)

\smallskip
\textit{Lecture on October 27, 2025, Slides 18--19, 24--28.}  

\end{enumerate}

% ------------------------------------------------------------------------------------------------------------------------------------
  
\item Development of Atomic Theory

\begin{enumerate}[label=\arabic*.]
\setcounter{enumii}{7} % Számláló beállítása

\item Describe the Rutherford scattering experiment. \hfill (5~points)

\smallskip
\textit{Lecture on October 20, 2025, Slides 5--7.}

\item What is the Zeeman effect? Describe the phenomenon. Under what conditions does it occur and how does it manifest itself? \hfill (5~points)

\smallskip
\textit{Lecture on October 20, 2025, Slides 16--17.}
  
\item Given the element with atomic number 11 in the periodic table, write down its electronic configuration using the standard atomic physics notation. \hfill (5~points)

\smallskip
\textit{Lecture on October 20, 2025, Slides 19--22.}
  
\end{enumerate}

\end{enumerate}

% ------------------------------------------------------------------------------------------------------------------------------------

Grades: 0–25~points: Fail (1), 26–30 points: Pass (2), 31–37 points: Satisfactory (3), 38–45~points: Good~(4), 46–50 points: Excellent (5).

\rightline{Dr.~Gabor FACSKO}
\rightline{\textit{facsko.gabor@uni-obuda.hu}}

\vfill

\end{document}
